\section{C. Smooth monodromy and the seven-dimensional invariant}

We distinguish the value of the invariant from formulas which would
compute it after a suitable coboundary has been constructed. In this
section $M_\Phi$ is a closed seven-manifold, whereas $W_8$ denotes a
compact oriented eight-manifold with boundary. The circle parameter
has positive orientation. On $M_\Phi$ we use the orientation
$d\theta\wedge o_{X_c}$; this is the boundary orientation of a
holomorphic family over a positively oriented disc. The statements
below do not assign a characteristic number to any unconstructed
filling in Part B.

\subsection{C.1. Clutching conventions and the mapping torus}

\begin{proposition}[C.1.1 (Known, with derivation): the mapping-torus comparison]
For an orientation-preserving smooth diffeomorphism
$\phi:S^6\to S^6$, put
\[
 M_\phi=(S^6\times[0,1])/((x,1)\sim(\phi(x),0)),\qquad
 \Sigma_\phi=D^7_+\cup_\phi(-D^7_-).
\]
The boundary identification uses the usual oriented parametrizations
of the two underlying copies of $S^6$; the minus sign reverses the
orientation of the second disc. Then
\[
 M_\phi\cong (S^1\times S^6)\mathbin\#\Sigma_\phi
\]
as oriented smooth manifolds. With the corresponding sphere-gluing
identification, $[\phi]\mapsto[\Sigma_\phi]$ identifies
$\pi_0\Diff^+(S^6)$ with $\Theta_7$. For either spin structure on
$M_\phi$,
\[
 \mu(M_\phi)=\mu(\Sigma_\phi).
\]
For the family in the prompt the gluing map is
$\Phi=F_{2,c_0-2}\circ H$, with $c(r)=c_0-2r$.
\end{proposition}

\begin{proof}
An ambient isotopy first arranges that $\phi$ is the identity on a
small closed disc: move the image of its centre back, correct the
oriented derivative through $\GL^+(6,\R)$, and use isotopy extension
on a smaller coordinate disc. Consequently its support lies in the
complementary disc $D^6$. Cut the mapping torus at its seam and take
the block $D^6\times[-\epsilon,\epsilon]$. The lower half of this
block has the old coordinates; on the upper half replace the
coordinates by $\phi^{-1}$. This identifies the reglued block with a
standard seven-disc. Its attaching map is $\phi$ on its upper
hemisphere and the identity on the lower hemisphere and side collar.
The upper hemisphere has boundary orientation $o_{D^6}$ for the
orientation $dt\wedge o_{D^6}$ on the block. Thus the replacement
has exactly the sign in $D^7_+\cup_\phi(-D^7_-)$.
Removing a small ball from the second disc in that sphere leaves the
same attaching construction; this proves the connected-sum formula.

Sphere gluing identifies pseudo-isotopy classes with homotopy spheres;
Cerf's pseudo-isotopy theorem changes pseudo-isotopy to isotopy for
the simply connected six-sphere. The exact primary statements are
Theorem 0 and Corollaries 1--2 of \citet[p.~8]{Cerf1970};
the sphere case is also explained in Theorem 1 and Corollary 1 of
\citet[p.~156]{Poenaru1971}. These statements compare the
sphere gluing and extension-over-a-disc constructions. This is the
high-dimensional theorem, not the separate statement $\Gamma_4=0$.
The oriented connected-sum group is Theorem 1.1 and Lemma 2.1 of
\citet[pp.~504--505]{KervaireMilnor1963}; their dimension table on
p.~504 gives order 28 in dimension seven.

The sphere summand is spin. The product $S^1\times S^6$ is spin and
has two spin structures; both bound $S^1\times D^7$ with the
corresponding choice on $S^1$. That coboundary has zero fourth
cohomology and zero signature, hence zero Eells--Kuiper invariant.
Boundary connected sum of this coboundary with a spin coboundary of
$\Sigma_\phi$ proves additivity directly: the fourth intersection
forms are an orthogonal sum and the Pontryagin evaluations add.
Thus the two boundary spin structures give the same value. This is
also the connected-sum property in \citet[Section 3]{EellsKuiper1962}.

Finally $e^{-\pi i(c_0-2r)}=e^{-\pi ic_0}e^{2\pi ir}$.
Parallel transport therefore reaches $X_{c_0-2}$ along the positive
$w$-circle; its endpoint is identified by $F_{2,c_0-2}$.
This confirms the sign and endpoint in R12. Two smooth connections
can be joined by a smooth path of horizontal distributions, producing
an isotopy of the resulting monodromies. Holomorphic automorphisms in
the connected group $\C^*$ also change the result only by isotopy.
\end{proof}

\begin{proposition}[C.1.2 (Known, with derivation): the positive generator]
The generator $1\in\Z/28$ is represented by the oriented sphere
bundle of a rank-four oriented real vector bundle on $S^4$ with
$e=1$ and $p_1=6$, oriented as the boundary of its disc bundle.
Its Eells--Kuiper invariant is $1/28$.
\end{proposition}
\begin{proof}
The clutching parameters $(h,j)=(2,-1)$ in the oriented
$S^3$-bundle construction give $e=h+j=1$ and $p_1=2(h-j)=6$;
we orient the base and fibre to give these signs. These are the
bundle conventions and computations of \citet[Section 6,
equations (12)--(13)]{EellsKuiper1962}. The Gysin sequence gives
the integral homology of $S^7$, and the homotopy sequence gives
simple connectivity. Thus the total space is a homotopy sphere.
Its disc bundle has intersection form $(1)$ on the zero section,
is spin because its second cohomology with $\Z/2$ coefficients
vanishes, and has $p_1^2=36$ and signature 1. Hence
\[
 \mu=\frac{36-4}{896}=\frac1{28}.
\]
The order-28 classification identifies this with the requested
positive generator. The integer and rational arithmetic is checked
in \path{code/c_exact.py}.
\end{proof}

\subsection{C.2. The corrected hypotheses for the characteristic formula}

\begin{theorem}[C.2.1 (Known, with derivation): Lead 6's formula]
Let $M=M_\phi$ have the orientation above. Let $W_8$ be a smooth
compact oriented spin$^c$ eight-manifold with $\partial W_8=M$.
Let $c\in H^2(W_8;\Z)$ be the first Chern class of its determinant
line, and choose a trivialization on the boundary which induces the
spin structure used on $M$. Write $p=p_1(TW_8)$, choose a relative
lift $\bar c$, and let $\bar p$ be the unique rational relative lift
of $p$. Then, with all products evaluated on $[W_8,M]$,
\begin{equation}\label{eq:c-spin-c}
 \mu(M)\equiv
 \frac{\bar p^{,2}-4\sigma(W_8)}{896}
 -\frac{\bar p\bar c^{,2}}{192}
 +\frac{\bar c^{,4}}{384}\pmod{\Z}.
\end{equation}
The value is independent of the lift $\bar c$. If $W_8$ is almost
complex, use its canonical spin$^c$ structure, for which
$c=c_1(TW_8)$ and $p=c_1^2-2c_2$. If $W_8$ is spin and its
determinant line is trivial, this reduces to
\[
 28\mu(M)=\frac{p_1^2-4\sigma(W_8)}{32}\pmod{28}.
\]
All the boundary cohomology and spin-structure conditions in this
statement hold for the mapping torus; the existence of the particular
geometric $W_8$ proposed in Part B is a separate assertion.
\end{theorem}

\begin{proof}
The mapping-torus exact sequence, or Proposition C.1.1, gives
$H^k(M;\Z)=0$ for $2\leq k\leq5$ and $H^1(M;\Z)=\Z$.
The same vanishing in degree two holds with $\Z/2$ coefficients,
so $M$ is spin. Its spin structures form a torsor for
$H^1(M;\Z/2)=\Z/2$. The exact sequence of the pair gives an
isomorphism
$H^4(W_8,M;\Q)\to H^4(W_8;\Q)$ and a surjection
$H^2(W_8,M;\Z)\to H^2(W_8;\Z)$, proving the assertions
about lifts and determinant trivializations.

If two lifts of $c$ differ by $\delta\alpha$ with
$\alpha\in H^1(M;\Z)$, the relative cup-product identity gives
$\delta\alpha\smile b=\delta(\alpha\smile i^*b)$.
The restrictions of $c$ and $p$ vanish on $M$. Every term in the
change of $\bar p\bar c^2$ or $\bar c^4$ is therefore zero;
in particular $(\delta\alpha)^2=0$. This proves lift independence.
There is no assertion here that the relative class $\bar c$ itself
is unique.

For a closed spin$^c$ eight-manifold the degree-eight part of the
Dirac index density is
\[
 [e^{c/2}\widehat A]_8
 =\frac{7p_1^2-4p_2}{5760}
       -\frac{p_1c^2}{192}+\frac{c^4}{384}.
\]
Eliminate $p_2$ using $\sigma=(7p_2-p_1^2)/45$.
The result is exactly the right side of
\eqref{eq:c-spin-c}. To justify its use on a boundary, glue $W_8$
to the negative of a spin coboundary inducing the same boundary spin
structure; extend the determinant line by its chosen trivialization.
The glued manifold is spin$^c$, so its Dirac index is an integer.
Novikov additivity and the relative-product comparison identify that
integer with the difference between \eqref{eq:c-spin-c} and the
spin formula for $M$. The spin formula is
\citet[Section 6, equation (11)]{EellsKuiper1962}.
The spin$^c$ version is also explicitly recorded, with normalization
$28\mu$ and characteristic class
$\widehat p=(p_1+c^2)/2$, by
\citet[Section 2.2, equation (2.1)]{CrowleyNordstrom2020}.
Substituting this $\widehat p$ gives the same coefficients; both
polynomial identities are verified exactly in \path{code/c_exact.py}.

A spin coboundary can be supplied using the connected-sum description:
bound the product by $S^1\times D^7$ and bound the sphere summand
by its spin coboundary. Thus no inference from simple connectivity of
$M$ is made; indeed $\pi_1(M)=\Z$. The original Eells--Kuiper
paper's unstructured formulation has an extra condition on
$H^1(W;\Z/2)\to H^1(M;\Z/2)$ to ensure spin structures glue.
Our explicit requirement that the boundary spin structures match
supplies precisely that gluing datum. By Proposition C.1.1 the two
possible spin structures on this particular $M$ have equal invariant.
Finally the splitting principle gives
$p_1(TW_8)=c_1^2-2c_2$ for an almost complex tangent bundle.
\end{proof}

\subsection{C.2 supplement. A computed coboundary with a different boundary}

\begin{proposition}[C.2.2 (Proved): the untwisted intersection form]
Let $W_A$ be a sufficiently small compact tubular representative of the
smooth total space $\mathcal X^+$ of A.2.2 over a closed $u$-disc.
Its boundary is the mapping torus of Engel's \emph{untwisted} family
$X(u)$. It is not the six-sphere mapping torus $M_\Phi$.
Write $\mathcal L=H^2(S;\Z)$, $F=c_1\mathcal O_S(S_\infty)$,
$m=c_1(L)$, and $v=6P-O-3F$ in the conventions of A.4.1.
Let $H$ be the infinity-section class on $R$ and $A$ the exceptional
Cartier divisor pulled back to $R^+$. Then
\[
 H_4(W_A;\Z)=
 \frac{\mathcal L\oplus\Z H\oplus\Z A}{\Z(m-A)}.
\]
For $\alpha=(a,h,b)\in H^2(R^+;\Z)$ set
\[
 d(\alpha)=(I-T^*)a+h(m-F)+bF.
\]
The middle intersection form of $W_A$ is
\[
 B(\alpha,\beta)=-d(\alpha)\cdot d(\beta).
\]
Its rank is two, its signature is $+2$, and its two nonzero Smith
factors are $1,11$. On the image modulo its radical it is represented,
in the ordered basis $(v,m)$, by
\[
 \begin{pmatrix}6&1\\1&2\end{pmatrix}.
\]
\end{proposition}
\begin{proof}
The local semistable charts $xy=u$ give a tubular representative which
retracts to $Q^+$; its boundary is the nearby-fibre mapping torus. One
may equivalently choose a regular neighbourhood of $Q^+$ in the smooth
total space and round its boundary. The fibrewise plumbing identifies
that boundary with the inverse image of a small positive $u$-circle.
The normalization mapping cylinder and $H_3(S)=0$ identify its fourth
homology with the cokernel of
$H_4(S)\to H_4(R^+)$. In the divisor coordinates of A.4.2 the two
sections are $D'_0=H+m-A$ and $D_\infty=H$. Their difference is
$m-A$, proving the group presentation.

The normalization is an immersed complex hypersurface in $W_A$.
Its normal line has class $-D'_0-D_\infty$, because the whole central
divisor is the divisor of the base coordinate $u$. Put
$r_0\alpha=a+bF$ and $r_\infty\alpha=a-hm$; these are the
restrictions to the two sections. The self-intersection of two
four-cycles pushed forward from the normalization contains the normal
Euler term and the intersections between distinct local branches.
All transverse complex branch intersections have positive sign.
Consequently it is
\begin{align*}
 &-r_0\alpha\cdot r_0\beta-r_\infty\alpha\cdot r_\infty\beta\\
 &\quad+r_0\alpha\cdot T^*r_\infty\beta
       +T^*r_\infty\alpha\cdot r_0\beta\\
 &=-\bigl(r_0\alpha-T^*r_\infty\alpha\bigr)
       \cdot\bigl(r_0\beta-T^*r_\infty\beta\bigr).
\end{align*}
This is an intersection-of-cycles calculation in the actual smooth
fourfold, rather than an assignment of a form to an unconstructed
central fibre. The isometry property of $T$ and $T^*m=m-F$ give
the stated expression for $d$. Also $d(m-A)=0$, so it descends.

The translation formula in A.4.1 shows that the image of $d$ is
$\langle F,v,m\rangle$. The surjectivity of the three intersection
coordinates proved there shows equality integrally. Its radical is
$\Z F$, since $F^2=Fv=Fm=0$ and
$v^2=-6$, $m^2=-2$, $vm=-1$. The remaining form is the displayed
positive matrix, of determinant eleven and with an entry equal to
one. Integral row and column reduction gives Smith factors $1,11$.
These computations, including the full eleven-by-eleven form, are
independently checked in \path{code/e_push_C_attempts.py}.
\end{proof}

\begin{proposition}[C.2.3 (Proved): a relative Pontryagin number for $W_A$]
In the coordinates of C.2.2, $p=p_1(TW_A)$ evaluates as
\[
 p(a,h,b)=-4(aF)-28h.
\]
It vanishes on the radical of the middle form and therefore has a
rational relative lift. The square of any such lift is
\[
 \bar p^{\,2}=\frac{4512}{11}.
\]
These numbers do not compute $\mu(M_\Phi)$. In particular the
canonical determinant class of $W_A$ restricts nontrivially to its
own boundary, so C.2.1 cannot be applied with that determinant.
\end{proposition}
\begin{proof}
Let $h:R^+\to R$ be the modification, $C$ its nodal elliptic-fibre
centre, and $c=c_1(R)=F+2H+m$. The centre is a regular embedding of
codimension two, also at its node: its ideal there is $(X,zv)$.
Its normal bundle has degree zero, and $c\cdot C=0$. The blow-up
and small resolution give
\[
 c_1(R^+)=h^*c-A,\qquad
 A\,h^*D\,h^*E=0,\quad
 A^2h^*D=-C\cdot D,\quad A^3=0.
\]
These are the regular-embedding blow-up identities. The subsequent
small resolution is crepant, and $A$ is the pullback of a Cartier
divisor, so they remain valid on $R^+$. The normal line to the
immersed normalization has class
$h^*N_0+A$, where $N_0=-(2H+m)$.

For completeness, the required second Chern pairings do not assume
that $R^+$ is K\"ahler. Analytic Hirzebruch--Riemann--Roch applies to
a compact complex manifold; see \citet[Theorem~4.3, pp.~564--565]{AtiyahSinger1968IIIE}.
The modification has $Rh_*\mathcal O_{R^+}=\mathcal O_R$.
For the blow-up this follows locally from its hypersurface description
in the projective line bundle and the vanishing of the direct images
of $\mathcal O_{\mathbb P^1}(-1)$; for the small resolution it follows
from the formal neighbourhood of its $(-1,-1)$ rational curve.
Thus $\chi(h^*D)=\chi(D)$ for every line bundle $D$ on $R$.
The linear coefficient in the Riemann--Roch polynomial for $kD$
and the preceding intersection identities give
\[
 c_2(R^+)h^*D=c_2(R)D+C\cdot D.
\]
Moreover $\chi(\mathcal O(A))=\chi(\mathcal O)=1$: on the
exceptional projective line bundle over $C$, the restriction is its
relative $\mathcal O(-1)$, which has zero derived direct image.
Pulling through the small resolution does not change this Euler
characteristic. Applying Riemann--Roch to $A$ yields $c_2(R^+)A=0$,
because $A^3$, $c_1(R^+)A^2$, and $c_1(R^+)^2A$ all vanish.

The real tangent bundle of $W_A$ restricted to the normalization is
the sum of its real tangent bundle and this normal line. Therefore
\[
 h^*p_1(R)+h^*N_0^2-4[C]
\]
has the same pairings on pulled-back divisors as $p_1(TW_A)$, and
the pairing with $A$ is zero. Here $[C]$ denotes its degree-four
cohomology class on $R$. Since $H^2=-mH$, we have
$(2H+m)^2=m^2=-2[\mathrm{pt}]$. The surface $S$ has
$p_1(S)=-24[\mathrm{pt}]$. The ruled-bundle tangent sequence thus
gives
\[
 p_1(R)+N_0^2=p_1(S)+2(2H+m)^2=-28[\mathrm{pt}],
\]
and $C\cdot(a+hH)=aF$. This proves the covector formula.

Modulo $\Z F$ the coordinate of $d(a,h,b)$ along $v$ is $aF$
and the coordinate along $m$ is $h$. Thus the covector on the
nondegenerate quotient is $(-4,-28)$. It annihilates the radical,
which is precisely the condition for a rational relative lift.
Its square is
\[
 (-4,-28)\frac1{11}
 \begin{pmatrix}2&-1\\-1&6\end{pmatrix}
 \binom{-4}{-28}=\frac{4512}{11}.
\]
Different relative lifts differ by a class pairing trivially with
absolute classes, so this value does not depend on the choice.

Finally the canonical first Chern class on the nearby untwisted
fibre is $f^*c_1\mathcal O_B(1)$. It is nonzero: the surjectivity of
$a\mapsto(aF,av,am)$ in A.4.1 supplies an integral two-cycle with
evaluation one, and the Mayer--Vietoris relations in A.4.2 do not
alter that evaluation. Hence no boundary trivialization of this
determinant exists. Even the formal expression
$(\bar p^{\,2}-4\sigma(W_A))/896=79/176$ is only an arithmetic
combination for this different boundary, not the invariant requested
in Part C. A cobordism implementing the finite logarithmic transforms
and the cusp closing map would be required to transport the result.
\end{proof}


\subsection{C.3. What an equivariant twist does and does not determine}

\begin{proposition}[C.3.1 (Proved): the circle action and local shear]
The effective circle subgroup $\rho=\R/\Z\subset\C/\Z$ of R6
acts freely outside its fixed double curve. That curve is a sphere,
and its two complex normal weights are $1,-1$. A monodromy
representative may be chosen $\rho$-equivariant. In marked real
coordinates on a regular torus fibre, its linear part is
$I-2\widehat\delta\otimes\gamma$; its affine correction is the
translation supplied by R2 and depends on the trivialization.
\end{proposition}

\begin{proof}
On a regular torus the circle acts by translation in the primitive
lattice direction $\widehat\delta$, and is therefore free.
At a multiple fibre, apply the invariant covector $\gamma$ to the
equation that a circle translation agrees with the $j$-th affine
deck generator to the power $a$. Since the deck translation has
$\gamma$-component $\pm1/m_j$, equality forces
$a\equiv0\pmod{m_j}$. Primitivity of $\widehat\delta$ then
forces the circle element to be the identity. This checks freeness
on the quotient using the actual twist vectors, including the root
centre; the argument is not just a statement about the product cover.

In a cusp toric chart, the relevant primitive edge directions are
$e_1,e_2,e_1-e_2$, with $e_2=\widehat\delta$.
Their nonzero determinants with $e_2$ have absolute value one.
Thus stabilizers off the edge in direction $e_2$ are trivial. Along
that edge the circle is fixed; the two normal characters in the
adjacent unimodular cones are $1$ and $-1$ because the base monomial
is invariant. The opposite-edge identification gives the double
curve with its two distinct triple-point endpoints, namely a
$\PP^1$. This verifies the toric step of R6 used here.

The circle acts on the entire family and commutes with the endpoint
isomorphism $F_2$. Averaging a Riemannian metric over this compact
circle makes the orthogonal horizontal distribution invariant.
Its parallel transport and hence the monodromy commute with $\rho$.
On a local marked torus chart,
$\Pi_{c-2}=\Pi_c(I-2\widehat\delta\otimes\gamma)$.
Inserting the identity complex-linear part and translation of R2
gives the asserted affine formula. A locally chosen lift of its
translation must not be multiplied by a real interpolation parameter
without checking its period cocycle; that operation need not descend
to the torus bundle.
\end{proof}

\begin{proposition}[C.3.2 (Proved): a global scalar gauge would be trivial]
Let a smooth circle act on a simply connected smooth manifold $Z$.
If $g:Z\to S^1$ is smooth and invariant under the circle, then
$z\mapsto g(z)\cdot z$ is smoothly isotopic to the identity.
The primitive coordinate $\theta$ detecting $\gamma$ on a regular
fibre of $X_c$ has no continuous extension to $X_c$.
\end{proposition}

\begin{proof}
Since $Z$ is simply connected, $g$ lifts to a smooth real function
$a$ with $g=e^{2\pi ia}$. Invariance implies
$a(\lambda z)-a(z)\in\Z$. This difference is continuous in the
connected circle and is zero at its identity, so it is zero
everywhere. The maps $z\mapsto e^{2\pi ita(z)}\cdot z$ are
diffeomorphisms with inverses obtained by replacing $t$ by $-t$.
They give the isotopy. In the present application a loop representing
$\widehat\gamma$ in a regular torus fibre has degree one under
$\theta$. If $\theta$ extended to the simply connected $S^6$,
that degree would be zero. The contradiction shows why the local
twist in Lead 7 is not a global scalar-gauge calculation.
\end{proof}

\begin{proposition}[C.3.3 (Known, with derivation): equivariance alone is insufficient]
All 28 oriented mapping classes of $S^6$ admit representatives
equivariant for one fixed semifree circle action with fixed set $S^2$.
Consequently the circle action, its fixed set, and its normal weights
do not determine $[\Phi]$.
\end{proposition}

\begin{proof}
Use $S^6\subset\operatorname{Im}\mathbb H\oplus\mathbb H$.
The equivariant Blakers--Massey construction gives a generator
$\sigma$ of $\pi_0\Diff^+(S^6)$ equivariant for simultaneous
quaternionic conjugation. The smooth reentrance statement is
\citet[Theorem 2.1]{DuranRigas2007}; Section 3.2 identifies this
representative and its powers with all 28 classes. Restrict
conjugation to $q=e^{i\theta/2}$, regarding $q$ and $-q$ as the
same effective circle element. Its fixed vector space consists of
$p\in\R i$ and $w\in\R+\R i$, a real three-space, so the
fixed set on the sphere is $S^2$. Outside it, at least one of the
two $(j,k)$-planes is nonzero, and an effective rotation fixing it
is the identity. Thus the action is semifree. All powers of
$\sigma$ commute with it. Conjugating the complex coordinate on
one of the two real normal planes writes the characters as
$1,-1$; this representation calculation is checked exactly in
\path{code/c_exact.py}. These are statements about smooth circle
actions and normal representations, not a comparison of their
complex structures with $X_c$.
\end{proof}

\begin{proposition}[C.3.4 (Proved): information lost under the proposed base change]
Under $w=r^6$ the smooth monodromy class becomes $6[\Phi]$.
Its value alone leaves two possibilities for $[\Phi]$ in $\Z/28$,
differing by 14. In particular a vanishing value after that base
change would not prove that $[\Phi]=0$.
\end{proposition}
\begin{proof}
The induced homomorphism on fundamental groups sends a positive
$r$-circle to six positive $w$-circles, so its transport is
$\Phi^6$ up to isotopy. The kernel of multiplication by six on
$\Z/28$ is $\{0,14\}$. The explicit preimages of all elements
of its image are checked in \path{code/c_exact.py}.
\end{proof}

\begin{proposition}[C.3.5 (Known formula, with proved comparison): the closing translation]
Fix the standard twists and all root and cusp markings of R2.
There is a coherent choice of the even-shift maps
\(\widetilde F_{h,c}:X_c\to X_{c+h}\), \(h\in2\mathbb Z\),
with the following explicit ordinary-chart translation. Choose global
roots \(F^3=t\), \(G^4=t-1\) on the uniformizing cover and put
\(\ell=(2\pi i)^{-1}\log(F/G)\pmod{\mathbb Z}\), \(e=(0,1)^t\).
Then the translation is
\begin{equation}\label{eq:C-explicit-closing}
 b_{h,c}=-\frac h{12}\Pi_c\widehat\gamma
              +\left(h\ell-\frac{h^2}{24}\right)e.
\end{equation}
For \(h=12\), it is the scalar multiplication
\begin{equation}\label{eq:C-twelve-scalar}
 K(t)=\frac{t^4}{(t-1)^3}
\end{equation}
on the ordinary \(\mathbb C^*\)-coordinate, after removing an integral
period. On the finite root charts its scalar representatives are
\(K(t)s_1^{-12}\) and \(K(t)s_2^{12}\), respectively, both units.
In a cusp coordinate \(q=t_c\), the scalar representative is
\(q^{-1}H(q)\), with \(H\) a holomorphic nonvanishing function.
The same map has a holomorphic \((\mathbb C^*)^2\)-unit representative
on the entire cusp filling. These two representatives differ by
exactly the \(\widehat\gamma\) cusp deck transformation.

For six positive \(w\)-loops the closing map is
\(\widetilde F_{12,c_0-12}\); thus the stated cusp exponent is \(-1\)
with the orientation convention of C.1.1. This computes the closing
cocycle, not the mapping class of its composition with parallel
transport. In particular it does not imply \(6[\Phi]=0\).
\end{proposition}
\begin{proof}
The raw formula without \(-h^2e/24\), its three extension checks,
and the coherent correction occur in
\cite[\texttt{p2:prop:global-even-shift} and the subsequent
orbit-groupoid calculations]{DGKXv1push}. We verify their exact use
here; the source is version~1, and version~2 omits these statements.
The global roots exist because the zero orders of \(t\) and \(t-1\)
on the simply connected uniformizing surface are divisible by three
and four. The local characters and their coprime orders give
\[
 \ell(g_1z)-\ell(z)=-\frac13,\qquad
 \ell(g_2z)-\ell(z)=\frac14,\qquad
 \ell(g_0z)-\ell(z)=\frac1{12}\pmod{\mathbb Z}.
\]
A change of logarithm adds an integral \(h\widehat\delta\) period,
so \eqref{eq:C-explicit-closing} is a well-defined torus translation.
Since \(\Pi_{c+h}=\Pi_c(I+hE)\), the two ordinary lattices agree;
the source-to-target lattice map of the identity linear part is
\(I-hE\).

For compatible logarithm lifts, covariance gives the affine defects
\[
 b_{h,c}(g_jz)-R_j(z)b_{h,c}(z)
       =\Pi_c(g_jz)\eta_j,\qquad
 \eta_1=\frac h2(\widehat u-\widehat w-\widehat\delta),\quad
 \eta_2=\frac h2(-\widehat w+\widehat\delta).
\]
Their \(\gamma\)-coordinates vanish, so the same right sides use
\(\Pi_{c+h}\). The vectors are integral for even \(h\). At the
cusp the defect is zero: \((A_0-I)\widehat\gamma=-\widehat\delta\)
cancels the increment of \(\ell\). These equations prove ordinary
and cusp descent with the stated orientation of the peripheral
loops.

For the finite comparison put \(w_j=(2\pi i)^{-1}\log s_j\).
The functions \(F/(s_1G)\) and \(Fs_2/G\) are invariant
holomorphic units on the respective full root discs. Choose
holomorphic logarithms \(\kappa_j\) of them. Then
\[
 \ell=\epsilon_jw_j+\kappa_j\pmod{\mathbb Z},\qquad
 \epsilon_1=1,\quad\epsilon_2=-1.
\]
The exact finite-to-ordinary comparison satisfies
\(\sigma_{j,c+h}-\sigma_{j,c}=h\epsilon_jw_je\).
Consequently the required finite-chart translation is
\[
 b^{(j)}_{h,c}
 =-\frac h{12}\Pi_c\widehat\gamma
       +\left(h\kappa_j-\frac{h^2}{24}\right)e.
\]
It is holomorphic on the full root disc and satisfies
\[
 b^{(j)}_{h,c}(g_jz)-R_j(z)b^{(j)}_{h,c}(z)
       -\frac{h\epsilon_j}{m_j}e
 =\Pi_c(g_jz)\eta_j.
\]
This is precisely equivariance for the source and target affine
finite actions, including the changed period column. The exact
identity
\(\sigma_{j,c+h}+b^{(j)}_{h,c}
=\sigma_{j,c}+b_{h,c}\) proves overlap agreement. Thus no real
interpolation of a multivalued period has been used.

At the cusp let \(q=e^{2\pi i s}\) and
\(C_{c,21}=\beta^{\rm part}+c+s\), which is holomorphic at
\(q=0\). The function
\[
 \widehat H=e^{2\pi i s/12}F/G
 \quad\hbox{satisfies}\quad
 \widehat H^{12}=q\frac{t^4}{(t-1)^3}.
\]
The right side is a unit because \(t\) has a simple pole at the
cusp. The character formula makes \(\widehat H\) single-valued;
it therefore extends as a nonvanishing holomorphic function. Write
\(\ell=-s/12+\kappa_0(q)\), with \(\kappa_0\) holomorphic.
Since
\(\Pi_c\widehat\gamma=(6\mu,-s+C_{c,21})^t\),
the logarithmic terms cancel in \eqref{eq:C-explicit-closing}:
\[
 \bar b_{h,c}(q)=
 \left(-\frac h2\mu(q),\,
 -\frac h{12}C_{c,21}(q)+h\kappa_0(q)-\frac{h^2}{24}\right)^t.
\]
Exponentiating gives a holomorphic unit in \((\mathbb C^*)^2\).
On every toric chart its action multiplies a character coordinate
by an integral character of this unit, so it and its inverse
extend across every double curve and both triple points. The
source and target cusp lattice multipliers differ by
\((1,e^{2\pi i h\lambda_1})=(1,1)\), for
\(\lambda\in\mathbb Z^2\). They are therefore identical. The
unit multiplication commutes with their toric shears and descends.
Exponentiation of the displayed formula gives exactly the ordinary
translation on the punctured cusp. We have now verified all
three comparisons and their inverses.

The uncorrected translations have composition defect
\(-hk e/12\). The constants in
\eqref{eq:C-explicit-closing} cancel it, since
\((h+k)^2-h^2-k^2=2hk\). Hence
\[
 \widetilde F_{k,c+h}\widetilde F_{h,c}
       =\widetilde F_{h+k,c},\qquad
 \widetilde F_{0,c}=\mathrm{id}.
\]
Any difference from the fixed R2 choice is a scalar automorphism,
which lies in its connected \(\mathbb C^*\). The choice therefore
has the same smooth monodromy class.

For \(h=12\), remove the actual period
\(-\Pi_c\widehat\gamma\) from \eqref{eq:C-explicit-closing}.
The remaining translation is \((12\ell-6)e\). Its exponential
is \((F/G)^{12}=K(t)\), proving
\eqref{eq:C-twelve-scalar}. The finite corrections divide it by
\(s_1^{12}\) and multiply it by \(s_2^{12}\). The original
uniformizers are retained: these are units because
\(t/s_1^3\) and \((t-1)/s_2^4\) are units. They need not be
replaced by one.

Put \(H(q)=qK(t(q))=e^{24\pi i\kappa_0(q)}\). The cusp unit
for \(h=12\) is exactly
\[
 U_{12,c}(q)=
 \left(e^{-12\pi i\mu(q)},\,
       e^{-2\pi i C_{c,21}(q)}H(q)\right).
\]
The \(\widehat\gamma\) deck transformation has multiplier
\[
 \left(e^{12\pi i\mu(q)},\,
       e^{2\pi i C_{c,21}(q)}q^{-1}\right),
\]
because \(B_0(1,0)^t=(0,-1)^t\). Multiplying these two
expressions gives \((1,q^{-1}H)=(1,K)\). This proves the
claimed comparison, including its sign, on the entire cusp
quotient. In particular the scalar representative's pole is not
an obstruction to the holomorphic map itself: its unit
representative already extends.

Finally \(w=e^{-\pi i c}\) makes six positive loops the path
\(c(r)=c_0-12r\); its endpoint is closed by the positive shift
\(h=12\). A quotient isotopy proving an order bound would still
have to patch the ordinary scalar description with the two-coordinate
cusp unit description. Scaling their additive logarithms separately
produces on their overlap a fractional
\(\widehat\gamma\)-period, which is not an integral period during
the interpolation. This patching has not been established. The
exact substitutions, coherent correction, finite defects, and
cusp comparison are checked in
\path{code/d_push_monodromy.py} and
\path{code/d_push_closing.py}; these are symbolic checks of the
identities proved above.
\end{proof}

\begin{proposition}[C.3.6 (Proved): independence of the equivariant lift]
Let \(X=X_c\) carry the effective semifree circle action of C.3.1,
with fixed surface \(F\cong S^2\). If two smooth equivariant
diffeomorphisms of \(X\) induce the same map of its orbit space,
they are smoothly isotopic. More generally, a smooth isotopy of
orbit maps which is supported in a compact subset of the free
orbit stratum lifts to an equivariant isotopy. Thus, whenever the
annular orbit twist has been constructed with a smooth equivariant
lift, its class in \(\Theta_7\) is independent of that lift.
This assertion does not construct an isotopy of the annular twist.
\end{proposition}
\begin{proof}
Their quotient is an equivariant orbit-preserving diffeomorphism
\(h\). On \(X\setminus F\) it has a unique expression
\(h(x)=g(x)\cdot x\), with \(g\) a smooth invariant circle-valued
function. We check its extension across \(F\) directly. Choose an
equivariant tubular slice and trivialize its normal bundle locally
over \(F\). Conjugating the negative-weight complex coordinate
identifies the normal action with scalar multiplication on
\(\mathbb C^2\). Since \(h\) preserves each orbit, it fixes the
base variable \(y\) and has normal part
\[
 (f_1(y,z),f_2(y,z))=g(y,z)(z_1,z_2),\qquad z\ne0.
\]
In particular \(z_2f_1=z_1f_2\), an identity of smooth functions
also at zero. In the homogeneous Taylor expansion in the four real
normal variables, each pair \((p_1,p_2)\) therefore satisfies
\(z_2p_1=z_1p_2\). In the complexified polynomial ring with
variables \(z_1,z_2,\bar z_1,\bar z_2\), this implies
\(p_1=z_1q\) and \(p_2=z_2q\): first set \(z_1=0\) and divide
by \(z_1\), then cancel \(z_1\). This coefficient argument works
with coefficients depending smoothly on \(y\). It produces a
common formal Taylor series for \(g\).

The formal series can be realized by a smooth complex function
\(G(y,z)\). Here is the parametric Borel construction needed for
this assertion. On a relatively compact base slice, write its
homogeneous terms as \(q_n(y,z)\), choose a smooth cutoff equal to
one near zero, and replace each term by
\(\chi(z/\epsilon_n)q_n(y,z)\). Choose
\(\epsilon_n\downarrow0\) so rapidly that the term and all
mixed base/normal derivatives of total order at most \(n/2\)
are bounded by \(2^{-n}\), away from the finitely many lower
degrees. This is possible because each such derivative retains a
positive power of \(\epsilon_n\). The resulting series converges
in every smooth seminorm, and its Taylor series is the prescribed
one: the cutoff is constant near zero and higher homogeneous
terms have zero lower jets. This proves the required local
parametric extension. The remainders
\(f_j-Gz_j\) are flat in the normal variables, locally uniformly
with all derivatives in \(y\). For \(z\ne0\),
\[
 g-G=\frac{\bar z_1(f_1-Gz_1)+\bar z_2(f_2-Gz_2)}{|z|^2}.
\]
A smooth function flat along \(z=0\), divided by \(|z|^2\),
extends smoothly and flatly by zero: the usual quotient derivative
formula and Taylor bounds of arbitrarily high order prove this for
every derivative. Thus \(g\) extends smoothly across the slice.
Its absolute value is one by continuity. Equivariance implies
invariance on the free part, hence on the slice by continuity.
The local extensions agree because they agree on the dense free
stratum. We obtain a global smooth invariant \(g:X\to S^1\).

Since \(X\) is simply connected, the covering
\(\mathbb R\to S^1\), \(a\mapsto e^{2\pi ia}\), gives a smooth
lift \(a\). The integer-valued difference
\(a(\lambda x)-a(x)\) is zero at \(\lambda=1\) and continuous
in the connected circle, so \(a\) is invariant. The maps
\(x\mapsto e^{2\pi isa(x)}\cdot x\), \(0\leq s\leq1\),
are diffeomorphisms with the explicit inverse obtained by changing
\(s\) to \(-s\). They give the required isotopy.

For the second assertion the free quotient is a principal circle
bundle. A connection lifts the compactly supported time-dependent
vector field of the orbit isotopy to an invariant vector field.
Its flow exists for the full finite time interval, is equivariant,
and extends by the identity near \(F\). The endpoint differs from
any other smooth equivariant lift of the same orbit map by the
orbit-preserving diffeomorphism just treated.

The general orbit-preserving statement is Haefliger and Salem
(1991, Theorem~3.1). We used the direct slice argument because the
non-equivariant Lemma~3.2 in that paper, as originally printed,
requires correction: complex conjugation on \(\mathbb C\)
preserves scalar-circle orbits but its quotient
\(\bar z/z\) has no continuous extension to zero. The theorem
with equivariance is unaffected; the corrected argument is also
explained in Karshon's research slides, which reproduce the original
statement. No conclusion about smooth maps of the topological orbit
space near \(F\) is assumed here.
\nocite{HaefligerSalem1991D,KarshonOrbitSlidesD}
\end{proof}


\subsection{C.4. Change of twists and real structures}

\begin{proposition}[C.4.1 (Proved): twist independence in the sphere cases]
For the fixed multiplicities $(3,4)$ and all admissible data in R5
with $|p|=1$, the monodromy class of the positive $w$-circle is the
same class in $\Theta_7$, with the complex orientation on each fibre.
\end{proposition}

\begin{proof}
The explicit sufficiency maps in E.1 depend holomorphically on a
small $c$-disc: the lattice defects and toric integral exponents are
constant, the periods and unit factors are holomorphic, and the
translation correction is solved in the fixed sheaf $V$ with
$H^1(B,V)=0$. This is the parameter-dependent assertion of E.1.3,
not just the fibrewise classification. For the second twist
coefficient $s_2$, the parameter identification is
$W=e^{-\pi i(c+s_2)}=(-1)^{s_2}w$.
The rotation of the punctured disc preserves its positive
fundamental-group generator. Local family identifications differ
by elements of $\Aut(X)=\C^*$, a connected group, so their
transition changes are smoothly isotopic to the identity. Following
these identifications around the circle conjugates the monodromy
in the oriented mapping-class group. Since $\Theta_7$ is abelian,
the resulting class is unchanged. If the lattice sign in R5 is
$-1$, the comparison map is nevertheless holomorphic and preserves
the complex orientation; no extra sign is introduced.
\end{proof}

\begin{proposition}[C.4.2 (Proved): conjugation gives no torsion constraint]
The antiholomorphic identification in R6 covers
$w\mapsto-\bar w$. At a fixed base point on the imaginary
$w$-axis it identifies the positive-loop monodromy with the
inverse-loop monodromy by an orientation-reversing fibre map.
Its action on $\Theta_7$ is negation, so the resulting equation is
the tautology $-[\Phi]=-[\Phi]$, not $2[\Phi]=0$.
\end{proposition}

\begin{proof}
From $c'=1-\bar c$ we obtain
$e^{-\pi ic'}=-\overline{e^{-\pi ic}}$.
Complex conjugation reverses the parameter circle. An
antiholomorphic map in complex dimension three reverses the real
orientation because its determinant has sign $(-1)^3$.
Conjugating a sphere clutching map by a reflection changes the
orientation of its glued homotopy sphere, hence negates its
$\Theta_7$ class. Equivalently the Eells--Kuiper invariant changes
sign. Applying this to the inverse loop proves the assertion.
This confirms the last sentence of Lead 7 without extracting a
spurious order-two restriction.
\end{proof}

\begin{problem}[C.5.1 (Open): the actual residue and its independent computation]
The value
\[
 [\Phi]\in\Z/28
\]
for the specified family has not been computed in this return.
Propositions C.1.1--C.4.2 do not supply it. To evaluate
\eqref{eq:c-spin-c} on the geometric route, one must construct the
specific resolution $W_8$ of a descended filling and compute its
relative intersection pairing, $c_1,c_2$, and signature. No such
intersection pairing is asserted here. On the equivariant route,
one must identify the global correction cocycle of the local
$\gamma$-shear with a class of equivariant sphere diffeomorphisms;
the regular-fibre shear and the fixed-point data alone do not do so.
The failed scalar-gauge reduction is excluded by C.3.2, and
C.3.3 explains why merely constructing an equivariant representative
cannot finish the calculation. A computation of $6[\Phi]$ would
still need to resolve the order-two ambiguity in C.3.4. Thus the
numeric requests in C.1--C.3 and the numeric part of C.4 remain
undecided, as does the full corrected-twist formula of Lead 7.
The characteristic formula in Lead 6 is established with its
hypotheses; its proposed geometric application is not.
\end{problem}
